\chapter{Vitamin B12 Deficiency}

\section*{Definition and Overview}
Vitamin B12 (also called cobalamin) is an essential nutrient that the body needs to make red blood cells and to keep the nervous system working normally. B12 deficiency means the body does not have enough of this vitamin, which can lead to anaemia and to damage of the nerves and spinal cord.

The body cannot make vitamin B12, so it must come from food. It is found naturally in meat, fish, eggs and dairy products. To be absorbed, vitamin B12 needs stomach acid and a protein made in the stomach called intrinsic factor. Deficiency can result from eating too little of the vitamin, from problems absorbing it, or from certain medicines. Left untreated, it can cause serious and sometimes irreversible problems, but it is usually easy to treat once diagnosed.

\section*{Epidemiology}
Vitamin B12 deficiency is common, especially in older adults, affecting perhaps one in ten people over 60 in some populations. It is more common in people who follow a strict vegan or vegetarian diet, because animal products are the main natural source of the vitamin.

People with certain conditions are at higher risk: pernicious anaemia, an autoimmune condition that prevents B12 absorption, is the most common cause in older people. Others at increased risk include those who have had weight-loss (bariatric) surgery or part of the stomach or small bowel removed, people with coeliac disease or Crohn's disease, and people taking long-term medicines such as metformin or acid-reducing drugs. In developing countries, dietary deficiency is a more important cause than in Western countries.

\section*{Aetiology and Pathophysiology}
B12 deficiency arises either because too little of the vitamin reaches the body or because the body cannot use it.

\begin{itemize}
    \item \textbf{Dietary deficiency:} a strict vegan or vegetarian diet without supplementation, or very poor overall nutrition
    \item \textbf{Pernicious anaemia:} the immune system attacks stomach cells that make intrinsic factor, so B12 cannot be absorbed
    \item \textbf{Gut disease:} coeliac disease, Crohn's disease or bacterial overgrowth in the small intestine interfere with absorption
    \item \textbf{Surgery:} removal of part of the stomach or ileum reduces absorption
    \item \textbf{Medicines:} long-term use of metformin or proton pump inhibitors can lower B12 levels
    \item \textbf{Other causes:} repeated use of nitrous oxide (laughing gas), and rare inherited conditions
\end{itemize}

The body needs B12 to make new cells and to maintain the protective coating (myelin) around nerves. Deficiency therefore slows blood cell production, causing a large-celled (megaloblastic) anaemia, and damages nerves, causing a condition called subacute combined degeneration of the spinal cord.

\section*{Clinical Features}
Symptoms can develop slowly and are easy to miss, sometimes being mistaken for ageing or stress.

\begin{itemize}
    \item Fatigue, weakness and feeling faint
    \item Shortness of breath and palpitations with activity
    \item Pale or slightly yellow skin
    \item A sore, red or smooth tongue
    \item Numbness, tingling or a "pins and needles" feeling in the hands and feet
    \item Unsteadiness and poor balance or clumsiness
    \item Memory problems, difficulty concentrating and confusion
    \item Low mood, irritability or depression
    \item Reduced vision, in severe cases
\end{itemize}

\section*{Differential Diagnosis}
The symptoms of B12 deficiency are shared with many other conditions.

\begin{itemize}
    \item Iron deficiency anaemia, which can coexist with B12 deficiency
    \item Folate deficiency, which causes the same type of anaemia
    \item Other causes of peripheral neuropathy, such as diabetes or alcohol misuse
    \item Hypothyroidism, which causes fatigue and cognitive symptoms
    \item Dementia or cognitive decline in older people
    \item Multiple sclerosis, which can cause numbness, tingling and unsteadiness
    \item Depression and chronic fatigue syndrome
\end{itemize}

\section*{When to Seek Medical Care}
See a doctor if you have persistent fatigue, unusual numbness or tingling, a sore tongue, or problems with balance or memory, especially if you are in a higher-risk group. Early diagnosis prevents permanent nerve damage.

\textbf{Call 000 (or local emergency services) for:}
\begin{itemize}
    \item Severe breathlessness, chest pain or feeling about to collapse
    \item Sudden weakness, confusion or difficulty walking
    \item Rapidly progressing numbness or loss of vision
\end{itemize}

\section*{Diagnosis and Investigations}
B12 deficiency is diagnosed with blood tests, and the cause is then investigated.

\begin{itemize}
    \item \textbf{Vitamin B12 blood level:} the first test, though a borderline result needs further testing
    \item \textbf{Full blood count and blood film:} to look for the large red cells of megaloblastic anaemia
    \item \textbf{Methylmalonic acid (MMA) and homocysteine:} more sensitive markers that rise when B12 is low
    \item \textbf{Folate and iron levels:} because these deficiencies look similar and often coexist
    \item \textbf{Intrinsic factor antibody test:} to confirm pernicious anaemia
    \item \textbf{Further tests:} such as endoscopy or small bowel investigations, if another absorption problem is suspected
\end{itemize}

\section*{Management}
Treatment replaces the vitamin and addresses the cause, so it does not recur.

\begin{itemize}
    \item \textbf{B12 injections:} hydroxocobalamin injections, given frequently at first and then at intervals, usually for life in pernicious anaemia
    \item \textbf{High-dose oral B12 tablets:} an option for mild deficiency due to diet or certain absorption problems
    \item \textbf{Dietary changes:} adding meat, fish, eggs, dairy and fortified foods, or supplements for vegans
    \item \textbf{Treating the underlying cause:} for example, managing coeliac disease or reviewing medicines that lower B12
    \item \textbf{Monitoring:} repeat blood counts and B12 levels to confirm recovery
    \item \textbf{Follow-up testing:} checking for gastric cancer risk in people with pernicious anaemia
\end{itemize}

\section*{Complications}
\begin{itemize}
    \item Anaemia, which if severe can strain the heart and cause heart failure
    \item Nerve and spinal cord damage that may be permanent if treatment is delayed
    \item Falls and fractures in older people due to unsteadiness
    \item Increased levels of homocysteine, which is linked to heart and blood vessel disease
    \item Possible increased risk of stomach cancer in pernicious anaemia, which is why monitoring is advised
\end{itemize}

\section*{Prognosis and Outcomes}
The outlook is excellent when B12 deficiency is found early. Anaemia usually corrects within a few weeks of starting treatment, and most people feel better within days to weeks. Neurological symptoms such as tingling and unsteadiness improve more slowly, over months, and some may be permanent if the deficiency has been present for a long time. With lifelong treatment, people with pernicious anaemia can expect to lead normal, healthy lives.

\section*{Prevention and Self-Care}
\begin{itemize}
    \item Eat a diet with adequate B12 from meat, fish, eggs, dairy or fortified foods
    \item If you are vegan or vegetarian, take a B12 supplement and choose fortified products
    \item If you have had weight-loss or bowel surgery, follow your doctor's advice about lifelong supplementation
    \item Ask for B12 testing if you are on long-term metformin or acid-reducing medicines
    \item Do not self-treat long-term with B12 injections without a diagnosis, as this can mask other problems
    \item See a doctor early for persistent tiredness, tingling or balance problems
\end{itemize}

\section*{Sources and Further Reading}
The following reputable organisations provide reliable, current information on vitamin B12 deficiency.

\begin{itemize}
    \item National Health Service. \textit{Information on vitamin B12 and folate deficiency, causes and treatment.} https://www.nhs.uk/conditions/vitamin-b12-or-folate-deficiency-anaemia/
    \item Mayo Clinic. \textit{Overview of vitamin deficiency anaemia and B12 deficiency.} https://www.mayoclinic.org/diseases-conditions/vitamin-deficiency-anemia
    \item MedlinePlus. \textit{Health information on vitamin B12 and anaemia.} https://medlineplus.gov/vitaminb12.html
    \item National Institutes of Health Office of Dietary Supplements. \textit{Fact sheet on vitamin B12 for health professionals and consumers.} https://ods.od.nih.gov
    \item MSD Manual. \textit{Consumer information on vitamin B12 deficiency.} https://www.msdmanuals.com/home/disorders-of-nutrition
    \item Pernicious Anaemia Society. \textit{Information and support for people with pernicious anaemia and B12 deficiency.} https://pernicious-anaemia-society.org
\end{itemize}
